\chapter{Security Assessment}

\section{Assessment Method}
This chapter evaluates security posture strictly from visible implementation and scripts. The analysis is threat-oriented: what is protected, what remains exposed, and what controls should be prioritized next.

\section{Security Controls Implemented}
\subsection{Confidentiality and Integrity}
AES-256-GCM is used for packet payload protection. When used with unique nonces and validated tags, this offers strong confidentiality and tamper detection.

\subsection{Connection Authentication}
Shared-key HMAC token exchange prevents unauthenticated peers from entering packet-forwarding state, assuming key secrecy and clock skew within validation bounds.

\subsection{Input Guarding}
Framing code constrains packet lengths and catches incomplete reads; crypto code validates key sizes and ciphertext minimum length.

\section{Threat Model Fit}
\begin{table}[H]
\centering
\caption{Threat Coverage Mapping}
\begin{tabular}{p{0.36\textwidth}p{0.22\textwidth}p{0.30\textwidth}}
\toprule
\textbf{Threat} & \textbf{Coverage} & \textbf{Rationale} \\
\midrule
Passive packet sniffing & Strong & AES-GCM encrypts payload bytes \\
In-transit tampering & Strong & GCM authentication tag validation \\
Unauthorized peer connect & Moderate & Shared-key token auth only \\
Replay of data packets & Weak & No sequence window / anti-replay state \\
Key compromise blast radius & Weak & Static key reuse across sessions \\
Traffic analysis & Weak & Timing/size metadata visible \\
\bottomrule
\end{tabular}
\end{table}

\section{Critical Gaps}
\subsection{No Perfect Forward Secrecy}
A static symmetric key means compromise reveals all traffic that was recorded and encrypted under that key.

\subsection{No Session Key Derivation}
The same secret underpins both auth and content encryption pathways. Separation of concerns is recommended for future versions.

\subsection{Limited Replay Defenses}
Authentication tokens include timestamps, but packet replay within active sessions is not robustly prevented.

\subsection{Single-Factor Secret Model}
No certificate identity, no per-client credentials, and no revocation model are present.

\section{Script and Ops Security Considerations}
Operational scripts manipulate host routing and iptables. Incorrect cleanup or partial execution may leave persistent forwarding rules. This is an operational security concern, not a cryptographic flaw.

\section{Security Improvement Prioritization}
\begin{table}[H]
\centering
\caption{Prioritized Security Enhancements}
\begin{tabular}{p{0.12\textwidth}p{0.38\textwidth}p{0.36\textwidth}}
\toprule
\textbf{Priority} & \textbf{Enhancement} & \textbf{Expected Benefit} \\
\midrule
P1 & Ephemeral handshake + session key derivation & Forward secrecy and compartmentalization \\
P1 & Packet sequence IDs + replay window & In-session replay mitigation \\
P2 & Key rotation and rekey signaling & Reduced key lifetime risk \\
P2 & Separate keys for auth and encryption & Better key hygiene and auditability \\
P3 & Optional obfuscation/padding & Reduced traffic analysis leakage \\
\bottomrule
\end{tabular}
\end{table}

\section{Hardening Checklist}
\begin{enumerate}[leftmargin=2em]
\item Replace static shared-key model with authenticated key exchange.
\item Add protocol version byte and minimum accepted version policy.
\item Introduce structured error telemetry without leaking key material.
\item Add stricter resource caps and connection rate limiting on server.
\item Harden scripts for safe rollback and deterministic cleanup.
\end{enumerate}

\section{Residual Risk Statement}
Current implementation can protect confidentiality/integrity in controlled educational environments where key handling is trusted and adversarial sophistication is limited. It should not be treated as production-secure VPN software without major security upgrades listed above.

\section{Conclusion}
WhisperTunnel demonstrates solid foundational primitives but intentionally omits advanced controls. This makes it excellent for learning and prototyping, while clearly indicating the engineering distance to production-grade security assurances.
